\documentclass[11pt,letterpaper]{article}

\usepackage[top=1in, bottom=1in, left=1in, right=1in, headheight=14pt]{geometry}
\usepackage{fontspec}
\setmainfont{DejaVu Serif}
\setsansfont{DejaVu Sans}
\setmonofont{DejaVu Sans Mono}

\usepackage{xcolor}
\usepackage{titlesec}
\usepackage{fancyhdr}
\usepackage{enumitem}
\usepackage{hyperref}
\usepackage{booktabs}
\usepackage{tabularx}
\usepackage{longtable}
\usepackage{colortbl}
\usepackage{multirow}
\usepackage{array}
\usepackage{parskip}
\usepackage{tcolorbox}
\usepackage{graphicx}
\usepackage{lastpage}

% === COLOR SCHEME: Ecology/Nature ===
\definecolor{primary}{HTML}{1B5E20}       % Forest green — sections
\definecolor{secondary}{HTML}{1565C0}     % River blue — subsections
\definecolor{tertiary}{HTML}{4E342E}      % Earth brown — subsubsections
\definecolor{accent}{HTML}{2E7D32}        % Highlights
\definecolor{lightprimary}{HTML}{E8F5E9}  % Quote boxes
\definecolor{lightsecondary}{HTML}{E3F2FD}
\definecolor{lighttertiary}{HTML}{EFEBE9}
\definecolor{rowgray}{HTML}{F5F5F5}       % Alternating table rows
\definecolor{bordergray}{HTML}{BDBDBD}    % Rules
\definecolor{subtlegray}{HTML}{757575}    % Footer text
\definecolor{darktext}{HTML}{212121}      % Body text

% === SECTION FORMATTING ===
\titleformat{\section}
  {\Large\bfseries\sffamily\color{primary}}
  {\thesection}{1em}{}
  [\vspace{-0.5em}\textcolor{bordergray}{\rule{\textwidth}{0.4pt}}]

\titleformat{\subsection}
  {\large\bfseries\sffamily\color{secondary}}
  {\thesubsection}{1em}{}

\titleformat{\subsubsection}
  {\normalsize\bfseries\sffamily\color{tertiary}}
  {\thesubsubsection}{1em}{}

\titlespacing*{\section}{0pt}{1.5em}{0.8em}
\titlespacing*{\subsection}{0pt}{1.2em}{0.5em}
\titlespacing*{\subsubsection}{0pt}{0.8em}{0.3em}

% === CUSTOM ENVIRONMENTS ===
\newcommand{\stageheader}[2]{%
  \noindent\colorbox{#2}{\makebox[\textwidth][l]{%
    \textcolor{white}{\bfseries\sffamily\hspace{6pt}#1\hspace{6pt}}}}%
  \vspace{0.5em}\par%
}

\newtcolorbox{asciibox}{
  colback=rowgray, colframe=bordergray,
  arc=1pt, boxrule=0.5pt,
  left=6pt, right=6pt, top=4pt, bottom=4pt,
  fontupper=\small\ttfamily,
  breakable
}

\newtcolorbox{quotebox}{
  colback=lightprimary, colframe=primary,
  arc=2pt, boxrule=0.5pt,
  left=12pt, right=12pt, top=8pt, bottom=8pt,
  breakable
}

\newcommand{\metafield}[2]{\textbf{\sffamily #1} #2\par}
\newcolumntype{L}[1]{>{\raggedright\arraybackslash}p{#1}}
\newcolumntype{C}[1]{>{\centering\arraybackslash}p{#1}}
\newcommand{\tableheader}[1]{\textbf{\sffamily\textcolor{white}{#1}}}

% === HEADERS/FOOTERS ===
\pagestyle{fancy}
\fancyhf{}
\fancyhead[L]{\small\sffamily\textcolor{subtlegray}{Bee \& Pollinator Research}}
\fancyhead[R]{\small\sffamily\textcolor{subtlegray}{GSD Mission Package}}
\fancyfoot[C]{\small\sffamily\textcolor{subtlegray}{Page \thepage\ of \pageref{LastPage}}}
\renewcommand{\headrulewidth}{0.4pt}
\renewcommand{\footrulewidth}{0pt}

\hypersetup{
  colorlinks=true,
  linkcolor=secondary,
  urlcolor=secondary,
  pdfauthor={GSD Ecosystem},
  pdftitle={Feed the Bees -- Mission Package},
  pdfsubject={Vision to Mission Pipeline},
}

\begin{document}

% === TITLE PAGE ===
\thispagestyle{empty}
\begin{center}
\vspace*{3cm}

{\small\sffamily\textcolor{primary}{\textbf{GSD RESEARCH MISSION PACKAGE}}}

\vspace{0.3cm}
\textcolor{bordergray}{\rule{8cm}{0.4pt}}

\vspace{1.5cm}
{\fontsize{28}{34}\selectfont\bfseries\sffamily\textcolor{primary}{Feed the Bees\\[0.3em]A PNW Pollinator Guide}}

\vspace{0.8cm}
{\Large\sffamily\textcolor{secondary}{Pacific Northwest --- Indoor \& Outdoor Growing for Pollinators}}

\vspace{0.5cm}
{\large\sffamily\itshape\textcolor{tertiary}{A Deep Research Study}}

\vspace{3cm}
{\sffamily\textcolor{accent}{Vision $\rightarrow$ Research $\rightarrow$ Mission}}\\[0.3em]
{\small\sffamily\textcolor{accent}{Full Pipeline Execution}}

\vspace{3cm}
{\small\sffamily\textcolor{subtlegray}{Date: April 3, 2026  |  Status: Ready for GSD Orchestrator}}\\[0.3em]
{\small\sffamily\textcolor{subtlegray}{Activation Profile: Squadron  |  Estimated: 8 context windows across 2 sessions}}
\end{center}
\newpage

% === TABLE OF CONTENTS ===
\tableofcontents
\newpage

% ===================================================================
% STAGE 1 --- VISION DOCUMENT
% ===================================================================
\stageheader{STAGE 1 --- VISION DOCUMENT}{primary}

\section{Feed the Bees --- Vision Guide}

\metafield{Date:}{April 3, 2026}
\metafield{Status:}{Research Complete / Ready for Mission Execution}
\metafield{Depends on:}{ECO cluster, MUK weather data, PNW bioregion baseline}
\metafield{Context:}{Artemis II mission, Rosetta Ecology cluster, Fox Companies habitat restoration vision}

\subsection{Vision}

The bees are disappearing. Not in some distant ecological abstraction, but here --- in Mukilteo, in Everett, in the gardens and empty lots and strips of grass between the sidewalk and the street. The question isn't whether we should do something about it. The question is: \textbf{what can one person plant this weekend that will feed a bee next month?}

This research mission answers that question for the Pacific Northwest --- specifically for the Puget Sound lowlands at 47--48\textdegree N latitude, where marine air, volcanic soil, and 37 inches of annual rainfall create one of the most productive growing environments on Earth. We cover indoor seed starting (because PNW spring arrives cold and wet), outdoor garden design (from a single pot to a full pollinator corridor), and the global context that makes this work urgent.

The framing is grassroots: anyone can do this. A bag of seeds, a pot, a window. The research is professional: every plant recommendation comes from university extension services, every population statistic from USDA or Xerces Society data, every growing guide from Master Gardener programs. Simple getting-started guides backed by real science.

\subsection{Problem Statement}

\begin{enumerate}
  \item \textbf{Pollinator habitat loss is accelerating.} Urban and suburban development replaces wildflower meadows with lawns, parking lots, and ornamental plants that provide no nectar or pollen. The Pacific Northwest has lost an estimated 90\% of its native prairie habitat.
  \item \textbf{Most people don't know what to plant.} Nurseries sell showy ornamentals bred for appearance, not pollinator value. Many popular garden plants are ``nectar deserts'' --- beautiful but nutritionally empty for bees. Worse, some are treated with systemic neonicotinoids.
  \item \textbf{Bloom gaps starve overwintering queens.} PNW native bees emerge as early as February and forage into October. Most gardens bloom June--August only. The critical early-spring and late-fall bloom gaps leave queens without food during the most vulnerable periods.
  \item \textbf{Indoor growing is underutilized for pollinator plants.} PNW gardeners start tomatoes indoors but rarely think to start pollinator wildflowers. Many native species benefit from cold stratification followed by indoor germination --- a natural fit for the PNW climate cycle.
  \item \textbf{Global bee decline requires local action at scale.} Colony Collapse Disorder, varroa mites, pesticide exposure, and climate disruption are global problems. But pollinator habitat is restored one garden at a time. Grassroots programs like Bee City USA and the Million Pollinator Garden Challenge demonstrate that distributed local action produces measurable regional impact.
\end{enumerate}

\subsection{Core Concept}

\textbf{One model:} A bloom-succession pollinator garden designed for the PNW climate, starting from seed indoors, transplanting outdoors, and providing continuous nectar/pollen from February through October --- feeding native bees through every critical lifecycle stage.

\begin{asciibox}
BLOOM SUCCESSION TIMELINE (PNW 47-48N)
=========================================
Feb  Mar  Apr  May  Jun  Jul  Aug  Sep  Oct
|    |    |    |    |    |    |    |    |
[CROCUS/WILLOW]          early queens
     [OREGON GRAPE/RED FLOWERING CURRANT]
          [CAMAS/LUPINE/PHACELIA]
               [NATIVE ROSES/PENSTEMON]
                    [LAVENDER/ECHINACEA]
                         [GOLDENROD/ASTER]
                              [SEDUM/WITCH HAZEL]
|----+----+----+----+----+----+----+----+---|
 Feb  Mar  Apr  May  Jun  Jul  Aug  Sep  Oct
=========================================
Goal: ZERO bloom gaps. Every week, something is flowering.
\end{asciibox}

\subsection{Research Modules}

\subsubsection{Module 1: PNW Pollinator Ecology}
Native bee diversity in western Washington. Species profiles, nesting requirements, foraging ranges, seasonal activity patterns. Threats and conservation status. The Puget Sound convergence zone as pollinator habitat.

\subsubsection{Module 2: Pollinator Garden Design}
Plant selection for PNW conditions. Bloom succession planning. Native vs. non-invasive non-native species. Container gardening, raised beds, bee lawns. Master Gardener and WSU Extension resources.

\subsubsection{Module 3: Indoor Growing for Pollinators}
Seed starting methods, cold stratification, timing for PNW transplant windows. Equipment (lights, heat mats, trays). Species-specific germination guides for pollinator plants.

\subsubsection{Module 4: Global Bee Conservation}
Colony Collapse Disorder research status. Global population trends (UN FAO, IPBES). Grassroots programs (Bee City USA, Pollinator Partnership, Million Pollinator Garden Challenge). Citizen science (Bumble Bee Watch, iNaturalist). Policy context.

\subsubsection{Module 5: Getting Started Guides}
Practical how-to documents for three entry points: (1) the windowsill garden, (2) the container patio, (3) the full pollinator bed. Each guide: seed list, timeline, care instructions, expected results. Written for absolute beginners.

\subsection{Chipset Configuration}

\begin{asciibox}
chipset:
  name: bee-pollinator-mission
  activation: squadron
  agents:
    FLIGHT:    Orchestrator (mission direction)
    PLAN:      Planner (wave decomposition)
    SCOUT:     Web Researcher (source gathering)
    EXEC-1:    Executor Track A (ecology + gardens)
    EXEC-2:    Executor Track B (conservation + guides)
    CRAFT-ECO: Ecology Specialist (species, habitat)
    CRAFT-GDN: Garden Specialist (plants, growing)
    VERIFY:    Fact-checker (source validation)
    INTEG:     Integration (cross-module links)
    CAPCOM:    Human interface
    BUDGET:    Token manager
    LOG:       Chronicle agent
\end{asciibox}

\subsection{Scope Boundaries}

\textbf{In Scope (v1.0):}
\begin{itemize}[nosep]
  \item PNW-specific plant recommendations (47--48\textdegree N, USDA zones 8a--8b)
  \item Native bee species profiles for western Washington
  \item Indoor seed starting guides for pollinator plants
  \item Bloom succession calendar (Feb--Oct)
  \item Container, raised bed, and in-ground garden designs
  \item Global conservation context and grassroots programs
  \item Getting-started guides for three experience levels
\end{itemize}

\textbf{Out of Scope:}
\begin{itemize}[nosep]
  \item Honeybee management/beekeeping (separate domain)
  \item Commercial agriculture pollination services
  \item Pesticide regulation advocacy (present data, not policy positions)
  \item Non-PNW climate zones (future missions)
  \item Butterfly-specific plants (overlap exists but distinct mission)
\end{itemize}

\subsection{Success Criteria}

\begin{enumerate}
  \item Bloom succession calendar covers every week from February through October with at least one species flowering
  \item At least 20 PNW-native pollinator plant species individually profiled with growing instructions
  \item At least 10 PNW native bee species individually profiled with habitat requirements
  \item Indoor seed starting guide includes timing, method, and expected germination for at least 12 species
  \item Container garden guide requires no more than \$50 in materials and fits a 4-foot balcony
  \item All plant recommendations sourced from WSU Extension, OSU Extension, or Xerces Society
  \item All bee population data sourced from USDA, Xerces Society, or peer-reviewed journals
  \item At least 3 grassroots programs profiled with participation instructions
  \item Getting-started guides are actionable by someone who has never gardened before
  \item Every species recommendation includes bloom time, sun requirements, water needs, and pollinator value
  \item Connection to at least 2 other Rosetta cluster projects (ECO, MUK, CWC)
  \item Through-line connects to Fox Companies habitat restoration vision
\end{enumerate}

\subsection{Relationship to Other Documents}

\begin{center}
\rowcolors{2}{rowgray}{white}
\begin{tabularx}{\textwidth}{L{4cm}X}
\toprule
\rowcolor{primary}
\tableheader{Document} & \tableheader{Relationship} \\
\midrule
ECO cluster (26 files) & Parent ecology project. Bee research becomes a module within ECO. \\
MUK weather data & KPAE microclimate data informs planting zone recommendations. \\
CWC culinary & Herb gardens overlap: lavender, thyme, rosemary are both culinary and pollinator plants. \\
Forest sim (vehicle.js) & Pollinator behavior can be modeled in the existing ecosystem simulation. \\
Fox Companies vision & Habitat restoration is a core pillar. This mission provides the research foundation. \\
\bottomrule
\end{tabularx}
\end{center}

\subsection{The Through-Line}

The bees don't need a plan. They need flowers. Every garden planted, every lawn converted, every window box filled with native wildflowers is a node in a distributed pollination network that no central authority could design or manage. This is the same architecture we build in software: distributed, resilient, locally autonomous, globally beneficial. The co-op network model --- food, energy, communication --- extends naturally to habitat. Feed the bees and you feed the system that feeds you.

\textit{Planting a flower is the smallest possible unit of habitat restoration. It is also the most scalable.}

\newpage

% ===================================================================
% STAGE 2 --- RESEARCH REFERENCE
% ===================================================================
\stageheader{STAGE 2 --- RESEARCH REFERENCE}{secondary}

\section{Feed the Bees --- Research Reference}

\subsection{How to Use This Document}

This reference compiles findings from government agencies, university extension services, and peer-reviewed research. All numerical claims are attributed to specific sources. Use this as the evidence base for mission execution.

\textbf{Key source organizations:}
\begin{itemize}[nosep]
  \item USDA Agricultural Research Service --- pollinator research, colony loss surveys
  \item Xerces Society for Invertebrate Conservation --- native bee guides, habitat restoration
  \item WSU Extension --- PNW-specific growing guides, Master Gardener resources
  \item OSU Extension --- Pacific Northwest pollinator plant lists
  \item UN FAO --- global pollination economics, biodiversity assessments
  \item IPBES --- Global Assessment Report on Pollinators
\end{itemize}

% === RESEARCH DATA PLACEHOLDER ===
% The following sections will be populated from research agent findings.
% Agents: bee-ecology, bee-gardens, bee-global

\subsection{Module 1: PNW Pollinator Ecology}

\subsubsection{Species Diversity}

Washington State hosts approximately \textbf{600 native bee species} across six families (WSU Extension). Combined with Oregon, over 900 sensitive species are documented (USDA Forest Service). The Pacific Northwest supports $\sim$30 bumble bee species, 200+ mason and leafcutter bees, and 75+ mason bee species in the genus \textit{Osmia} alone.

\begin{center}
\rowcolors{2}{rowgray}{white}
\begin{tabularx}{\textwidth}{L{3.5cm}C{2cm}X}
\toprule
\rowcolor{primary}
\tableheader{Family} & \tableheader{PNW Species} & \tableheader{Key Genera \& Notes} \\
\midrule
Halictidae (sweat bees) & Most abundant & \textit{Lasioglossum} most species-rich genus (21+ in Puget Sound). Ground-nesting. \\
Apidae (bumble bees) & $\sim$30 \textit{Bombus} & Only native social bees in PNW. \textit{Melissodes} (long-horned) on Asteraceae. \\
Megachilidae (mason/leafcutter) & 200+ & \textit{Osmia lignaria} (blue orchard bee) pollinates 95\% of flowers visited. \\
Andrenidae (mining bees) & Diverse & \textit{Andrena} ubiquitous. Multiple new county records (WA Bee Atlas 2025--2026). \\
Colletidae (plasterer bees) & $\sim$33 & \textit{Hylaeus} nests in stems and wood cavities. \\
Melittidae (oil-collecting) & Rare & Specialist species, uncommon but present. \\
\bottomrule
\end{tabularx}
\end{center}

\subsubsection{Nesting Requirements}

\textbf{Ground-nesting ($\sim$70\% of native bee species):} Excavate tunnels 6--36+ inches deep in bare, well-drained soil. South-facing slopes preferred. Mowing, tilling, and compaction destroy nests. Includes \textit{Andrena}, \textit{Lasioglossum}, \textit{Halictus}, \textit{Colletes}, and \textit{Bombus} queens (abandoned rodent burrows).

\textbf{Cavity-nesting ($\sim$30\%):} Nest in beetle galleries, hollow stems, snail shells. Mason bees use mud; leafcutter bees cut leaf pieces. Key PNW pithy-stemmed plants for nesting: elderberry, oceanspray, mock orange, thimbleberry, snowberry, Pacific ninebark, roses, and raspberry/blackberry canes. Leave dead wood and stems standing through winter.

\subsubsection{Seasonal Bloom Calendar (47--48\textdegree N)}

\textbf{Early Spring (Feb--Apr):} Red flowering currant (\textit{Ribes sanguineum}) --- critical early food for emerging bumble bee queens. Willow (\textit{Salix} spp.) --- earliest pollen source. Oregon grape (\textit{Mahonia aquifolium}). Indian plum (\textit{Oemleria cerasiformis}).

\textbf{Spring--Early Summer (Apr--Jun):} Camas (\textit{Camassia quamash}) --- historically critical PNW species. Broadleaf lupine (\textit{Lupinus latifolius}) --- buzz-pollinated by bumble bees. Mock orange (\textit{Philadelphus lewisii}) --- WA state flower. Salal (\textit{Gaultheria shallon}).

\textbf{Summer (Jun--Aug):} Oceanspray (\textit{Holodiscus discolor}) --- also provides pithy nesting stems. Pacific ninebark. Oregon sunshine (\textit{Eriophyllum lanatum}). Nootka rose (\textit{Rosa nutkana}).

\textbf{Late Summer--Fall (Aug--Oct):} Douglas aster (\textit{Symphyotrichum subspicatum}) --- \textbf{\#1 ranked native plant for bee visitation} in OSU research. Canada goldenrod (\textit{Solidago canadensis}) --- critical pre-hibernation food for bumble bee queens. Puget Sound gumweed (\textit{Grindelia integrifolia}).

\textit{OSU finding: Native straight species attract significantly more pollinators than cultivated varieties (cultivars). Source: OSU Extension EM-9363.}

\subsubsection{Threats}

\textbf{Habitat loss (\#1 threat):} Puget Sound lowlands have lost 90\%+ of prairie/savanna habitat since European settlement. Invasive plants (Scotch broom, Himalayan blackberry, English ivy) displace native forage.

\textbf{Pathogen spillover:} Commercially reared bumble bees transmit \textit{Nosema bombi} and \textit{Crithidia bombi} to wild populations. Possibly the most significant factor in \textit{B. occidentalis} collapse ($>$90\% decline). Prevalence increases near greenhouse operations.

\textbf{Climate:} Flowering has advanced 8+ days since 1951. WSU 2024 research projects spring colony collapse from warmer falls --- outdoor colonies in eastern WA projected below collapse threshold (9,000 adults) by 2050.

\subsubsection{Convergence Zone as Habitat}

Sugden et al.\ (2025) documented \textbf{118 bee species in just 3 marginal urban sites} in the Puget Sound region --- significantly more than diversified farms. Power line corridors and airport perimeters function as unexpected biodiversity hotspots: regularly disturbed, open, sun-exposed, with bare ground for nesting. The mild maritime climate extends the active season (Feb--Oct), but heavy urbanization and invasive species remain major challenges.

\subsection{Module 2: Pollinator Garden Design}

\subsubsection{Top Native Plants for PNW Pollinator Gardens}

All recommendations sourced from WSU Extension, OSU Extension, Xerces Society Maritime NW plant list, and USDA NRCS.

\textbf{Native Wildflowers (selected highlights):}

\begin{center}
\rowcolors{2}{rowgray}{white}
\begin{tabularx}{\textwidth}{L{3cm}L{3cm}C{1.5cm}X}
\toprule
\rowcolor{primary}
\tableheader{Common Name} & \tableheader{Latin Name} & \tableheader{Bloom} & \tableheader{Pollinator Value} \\
\midrule
Common Camas & \textit{Camassia quamash} & Apr--May & Bumble bees, mason bees \\
Oregon Sunshine & \textit{Eriophyllum lanatum} & May--Jul & 20+ bee genera documented \\
Broadleaf Lupine & \textit{Lupinus latifolius} & May--Jul & Buzz-pollinated by bumble bees \\
Cascade Penstemon & \textit{Penstemon serrulatus} & Jun--Aug & Bumble bees, mason bees, hummingbirds \\
Blanketflower & \textit{Gaillardia aristata} & Jun--Sep & Long bloom; bees, butterflies \\
Douglas Aster & \textit{S. subspicatum} & Aug--Oct & \textbf{\#1 bee visitation (OSU)} \\
Canada Goldenrod & \textit{Solidago canadensis} & Aug--Oct & Critical pre-hibernation food \\
Showy Milkweed & \textit{Asclepias speciosa} & May--Aug & Monarch host + bumble bees \\
\bottomrule
\end{tabularx}
\end{center}

\textbf{Native Shrubs:} Red flowering currant (\textit{Ribes sanguineum}, Mar--Apr) --- first major spring forage for bumble bee queens. Oregon grape (\textit{Mahonia aquifolium}, Mar--May). Indian plum (\textit{Oemleria cerasiformis}, Feb--Mar) --- earliest PNW bloom. Mock orange (\textit{Philadelphus lewisii}, May--Jun) --- WA state flower. Oceanspray (\textit{Holodiscus discolor}, Jun--Jul) --- also provides pithy nesting stems.

\textbf{Non-Native, Non-Invasive (Extension-recommended):} Lavender (\textit{Lavandula angustifolia}), catmint (\textit{Nepeta} spp.), purple coneflower (\textit{Echinacea purpurea}), anise hyssop (\textit{Agastache} spp.), sedum (\textit{Sedum spectabile}, Aug--Oct --- late-season critical), hardy geranium `Rozanne' (May--Oct --- longest bloom of any perennial).

\subsubsection{Bloom Succession Strategy}

The goal: \textbf{at least 3 species flowering in every month from March through October.} Design principle from WSU Extension: ``Plant at least three species per bloom window. Overlap is intentional and beneficial.''

\begin{center}
\rowcolors{2}{rowgray}{white}
\begin{tabularx}{\textwidth}{C{1.5cm}C{1.5cm}X}
\toprule
\rowcolor{primary}
\tableheader{Month} & \tableheader{Species} & \tableheader{Key Plants} \\
\midrule
Mar & 6 & Indian plum, Oregon grape, kinnikinnick, red currant, pasque flower, bleeding heart \\
Apr & 9 & Camas, columbine, serviceberry, Oregon iris, elderberry, twinberry, balsamroot, strawberry, chocolate lily \\
May & 12 & Oregon sunshine, mock orange, ninebark, Nootka rose, lupine, milkweed, salal, catmint, lavender begins \\
Jun & 8 & Cascade penstemon, tiger lily, oceanspray, hardhack, blanketflower, echinacea, globe thistle \\
Jul & 6+ & Continuing June species + agastache, sedum begins, Richardson's penstemon \\
Aug & 6 & Douglas aster, goldenrod, pearly everlasting, sedum, gumweed, asters \\
Sep & 4 & Douglas aster, goldenrod, sedum, late asters \\
Oct & 3 & Sedum, witch hazel, late goldenrod \\
\bottomrule
\end{tabularx}
\end{center}

\subsubsection{Bee Lawns}

WSU and UMN Extension recommend replacing portions of traditional grass with low-growing flowering plants that tolerate mowing at 3--4 inches:

\begin{itemize}[nosep]
  \item \textbf{Dutch white clover} (\textit{Trifolium repens}) --- nitrogen-fixing, self-seeding, 6--8 inch height
  \item \textbf{Self-heal} (\textit{Prunella vulgaris}) --- native PNW groundcover, purple flowers Jun--Sep
  \item \textbf{Creeping thyme} (\textit{Thymus serpyllum}) --- fragrant, drought-tolerant, foot-traffic resistant
  \item \textbf{English daisy} (\textit{Bellis perennis}) --- naturalized in PNW, blooms Mar--Oct
\end{itemize}

Method: Overseed into existing lawn at 2--3 oz/1000 sq ft after core aeration. Reduce mowing frequency to every 2--3 weeks. Stop herbicide use. Results visible in 2--3 months.

\subsubsection{Container \& Small-Space Gardens}

For balconies, patios, and windowsills. All recommendations from WSU and OSU Extension container gardening guides:

\begin{itemize}[nosep]
  \item Minimum pot size: 12 inches diameter, 10 inches deep
  \item Essential trio for a single large container: lavender + catmint + sedum (covers Jun--Oct)
  \item Window box: pansies (early) $\rightarrow$ herbs (mid) $\rightarrow$ asters (late)
  \item Hanging basket: trailing nasturtiums + creeping thyme
  \item Use potting mix (not garden soil). Water consistently. Full sun preferred.
\end{itemize}

\subsection{Module 3: Global Bee Conservation}

\subsubsection{Colony Collapse Disorder}

Colony Collapse Disorder (CCD) is defined by the EPA as the phenomenon where the majority of worker bees disappear, leaving behind a queen, food, and immature bees. First identified in winter 2006--2007, CCD has evolved from an acute crisis to a chronic condition.

\textbf{2025 USDA-ARS Breakthrough:} Dr.\ Judy Chen at the Bee Research Laboratory (Beltsville, MD) identified \textbf{viruses from amitraz-resistant Varroa mites} as the proximate cause of recent mass collapses. Deformed wing virus A/B and acute bee paralysis virus were found in all sampled bees from collapsed colonies. Amitraz resistance was detected in ``virtually all collected Varroa'' --- amitraz being the critical miticide used by most US beekeepers. \textit{Source: USDA ARS Research News, June 2, 2025.}

\textbf{Colony loss scale (2024--2025):} Over 60\% of commercial beekeeping colonies lost since prior summer (as of February 2025). 1.7 million colonies affected. Estimated financial impact: \$600 million. \textit{Source: USDA NASS Honey Bee Colonies Report, August 1, 2025.}

\subsubsection{Economic Value of Pollination}

\begin{center}
\rowcolors{2}{rowgray}{white}
\begin{tabularx}{\textwidth}{L{4.5cm}L{3cm}X}
\toprule
\rowcolor{secondary}
\tableheader{Metric} & \tableheader{Value} & \tableheader{Source} \\
\midrule
Global crop production dependent on pollinators & \$235--577 billion/year & IPBES 2016 \\
Wild flowering plants dependent on animal pollination & $\sim$90\% & IPBES 2016 \\
US pollination services value & \$20--29 billion/year & Cornell / USDA \\
Native pollinator services (US) & \$3 billion/year & Xerces Society \\
World food production dependent on bees & $\sim$33\% & UN FAO \\
\bottomrule
\end{tabularx}
\end{center}

\subsubsection{Grassroots Programs}

\textbf{Bee City USA} (Xerces Society, founded 2012): 300+ affiliate programs across cities, counties, colleges. Requirements include standing pollinator committee, native plant lists, IPM plan, and community events.

\textbf{Million Pollinator Garden Challenge} (National Pollinator Garden Network, 2015--2021): Achieved 1,040,000 gardens registered by December 2018, exceeding the million-garden goal. Partners included NWF, Pollinator Partnership, and American Seed Trade Association.

\textbf{Bee Better Certified} (Xerces Society, 2017): First third-party certification focused on pollinator conservation. Tens of thousands of certified acres. Retail presence at Kroger, Trader Joe's, Costco. Silk (Danone) committed to 100\% almond acreage applying Bee Better practices by 2025.

\subsubsection{Citizen Science}

\textbf{Great Sunflower Project} (Dr.\ Gretchen LeBuhn, SFSU, 2008): Over 100,000 members --- the largest pollinator-focused citizen science project. Participants count pollinators visiting plants in gardens, yards, and parks.

\textbf{Bumble Bee Watch} (Xerces Society + partners): $\sim$75,000 observations across 49 states. 2024: 7,800+ sightings in one year. Companion Bumble Bee Atlas program deploys trained community scientists for systematic surveys.

\textbf{iNaturalist}: 2024 Northeast US data: 31,953 pollinator observations, 1,136 species, 3,126 observers. Most observed: \textit{Bombus impatiens} (2,062 observations). Vermont: 2,500+ community scientists documenting 350+ wild bee species.

\subsubsection{Policy Context}

\textbf{EU Neonicotinoid Ban:} Full outdoor ban on three neonicotinoids (imidacloprid, clothianidin, thiamethoxam) effective May 2018. Thiacloprid withdrawn 2020. ECJ closed emergency authorization loophole January 2023.

\textbf{US Pollinator Health Task Force} (2014): Three goals --- reduce overwinter losses to 15\%, increase monarch population to 225M, restore 7M acres of habitat. EPA proposed cancellation of imidacloprid residential turf spray uses (2020).

\textbf{ESA Listings:} Rusty patched bumble bee (\textit{Bombus affinis}) listed Endangered. Franklin's bumble bee (\textit{Bombus franklini}) listed Endangered --- not seen since 2006. American bumble bee under ESA review.

\subsection{Safety and Sensitivity Considerations}

\begin{center}
\rowcolors{2}{rowgray}{white}
\begin{tabularx}{\textwidth}{L{3cm}XL{2cm}}
\toprule
\rowcolor{secondary}
\tableheader{Area} & \tableheader{Consideration} & \tableheader{Action} \\
\midrule
Pesticide data & Present EPA/USDA findings without advocacy & ANNOTATE \\
Indigenous plant knowledge & Credit specific tribal nations by name & ABSOLUTE \\
Endangered species & No specific nest/breeding site locations & ABSOLUTE \\
Invasive species & Distinguish ``non-native beneficial'' from ``invasive'' & GATE \\
Health claims & No medicinal plant claims without peer review & GATE \\
\bottomrule
\end{tabularx}
\end{center}

\newpage

% ===================================================================
% STAGE 3 --- MISSION PACKAGE
% ===================================================================
\stageheader{STAGE 3 --- MISSION PACKAGE}{tertiary}

\section{Milestone Specification}

\subsection{Mission Objective}

Produce a comprehensive, publication-quality research document on PNW pollinator gardening --- from native bee ecology through practical getting-started guides --- sourced entirely from government agencies, university extension services, and peer-reviewed research. The output becomes a module in the Rosetta Ecology cluster and a foundation for the Fox Companies habitat restoration vision.

\subsection{Architecture Overview}

\begin{asciibox}
WAVE EXECUTION PLAN
====================
Wave 0: Foundation     [HAIKU]     Sequential
  |- Shared schemas, source index, plant taxonomy template

Wave 1: Research       [SONNET+OPUS]  Parallel (2 tracks)
  |- Track A: PNW Ecology + Garden Design    [CRAFT-ECO + EXEC-1]
  |- Track B: Global Conservation + Guides   [CRAFT-GDN + EXEC-2]

Wave 2: Synthesis      [OPUS]      Sequential
  |- Cross-module integration, bloom calendar assembly
  |- Getting-started guide compilation

Wave 3: Publication    [SONNET]    Sequential
  |- HTML page generation, index linking
  |- Verification + fact-checking pass
\end{asciibox}

\subsection{Deliverables}

\begin{center}
\rowcolors{2}{rowgray}{white}
\begin{tabularx}{\textwidth}{C{0.5cm}L{4cm}XL{2cm}}
\toprule
\rowcolor{tertiary}
\tableheader{\#} & \tableheader{Deliverable} & \tableheader{Success Criteria} & \tableheader{Spec} \\
\midrule
1 & PNW Bee Species Guide & 10+ species profiled with habitat data & 01-ecology \\
2 & Bloom Succession Calendar & Every week Feb--Oct covered & 02-garden \\
3 & Plant Profile Database & 20+ species with full growing data & 02-garden \\
4 & Indoor Seed Starting Guide & 12+ species with germination methods & 03-indoor \\
5 & Global Conservation Overview & 3+ grassroots programs, UN/USDA data & 04-global \\
6 & Getting Started: Windowsill & Actionable for beginners, <\$20 & 05-guides \\
7 & Getting Started: Container & 4-foot balcony, <\$50 & 05-guides \\
8 & Getting Started: Garden Bed & Full pollinator bed design & 05-guides \\
9 & HTML Research Pages & Integrated into tibsfox.com/Research/ & 06-publish \\
\bottomrule
\end{tabularx}
\end{center}

\subsection{Component Breakdown}

\begin{center}
\rowcolors{2}{rowgray}{white}
\begin{tabularx}{\textwidth}{L{3.5cm}L{2.5cm}C{1.5cm}C{2cm}}
\toprule
\rowcolor{tertiary}
\tableheader{Component} & \tableheader{Dependencies} & \tableheader{Model} & \tableheader{Est. Tokens} \\
\midrule
00-shared-types & None & Haiku & 2K \\
01-pnw-ecology & 00-shared-types & Sonnet & 15K \\
02-garden-design & 00-shared-types & Opus & 20K \\
03-indoor-growing & 00, 02 & Sonnet & 12K \\
04-global-conservation & 00-shared-types & Sonnet & 15K \\
05-getting-started & 00, 02, 03 & Opus & 18K \\
06-publish-html & All above & Sonnet & 10K \\
07-verification & All above & Sonnet & 8K \\
\bottomrule
\end{tabularx}
\end{center}

\subsection{Model Assignment Rationale}

\begin{center}
\rowcolors{2}{rowgray}{white}
\begin{tabularx}{\textwidth}{C{1.5cm}L{3cm}X}
\toprule
\rowcolor{tertiary}
\tableheader{Model} & \tableheader{Components} & \tableheader{Rationale} \\
\midrule
Opus (25\%) & 02, 05 & Garden design requires judgment (plant compatibility, aesthetic, succession). Getting-started guides require empathy and clarity for beginners. \\
Sonnet (65\%) & 01, 03, 04, 06, 07 & Survey compilation, structured content, verification --- reliable pattern work. \\
Haiku (10\%) & 00 & Shared schemas, plant taxonomy template --- pure scaffold. \\
\bottomrule
\end{tabularx}
\end{center}

\subsection{Test Plan}

\begin{center}
\rowcolors{2}{rowgray}{white}
\begin{tabularx}{\textwidth}{L{3cm}C{1.5cm}C{1.5cm}C{2cm}}
\toprule
\rowcolor{tertiary}
\tableheader{Category} & \tableheader{Count} & \tableheader{Priority} & \tableheader{Failure Action} \\
\midrule
Safety-critical & 5 & Mandatory & BLOCK \\
Core functionality & 14 & Required & BLOCK \\
Integration & 8 & Required & BLOCK \\
Edge cases & 5 & Best-effort & LOG \\
\midrule
\textbf{Total} & \textbf{32} & & \\
\bottomrule
\end{tabularx}
\end{center}

\subsection{Mission Crew Manifest}

\begin{center}
\rowcolors{2}{rowgray}{white}
\begin{tabularx}{\textwidth}{L{2cm}L{3cm}X}
\toprule
\rowcolor{tertiary}
\tableheader{Role} & \tableheader{Agent} & \tableheader{Responsibility} \\
\midrule
FLIGHT & Orchestrator & Mission direction; go/no-go gates \\
PLAN & Planner & Wave decomposition; dependency management \\
SCOUT & Web Researcher & Source gathering; evidence compilation \\
EXEC-1 & Executor (Track A) & Ecology + garden design documents \\
EXEC-2 & Executor (Track B) & Conservation + getting-started guides \\
CRAFT-ECO & Ecology Specialist & Species profiles, habitat analysis \\
CRAFT-GDN & Garden Specialist & Plant selection, growing guides \\
VERIFY & Fact-checker & Source validation; accuracy checking \\
INTEG & Integration Agent & Cross-module bloom calendar assembly \\
CAPCOM & HITL Interface & Human approval at wave boundaries \\
BUDGET & Token Manager & Budget tracking; session planning \\
LOG & Chronicle Agent & Commit messages; milestone summaries \\
\bottomrule
\end{tabularx}
\end{center}

\subsection{Execution Summary}

\begin{center}
\rowcolors{2}{rowgray}{white}
\begin{tabularx}{\textwidth}{L{5cm}X}
\toprule
\rowcolor{tertiary}
\tableheader{Metric} & \tableheader{Value} \\
\midrule
Activation profile & Squadron (12 roles) \\
Total components & 8 (including shared types) \\
Parallel tracks & 2 (ecology+garden | conservation+guides) \\
Estimated context windows & 8 \\
Estimated sessions & 2 \\
Model split & 25\% Opus / 65\% Sonnet / 10\% Haiku \\
Total estimated tokens & $\sim$100K output \\
Test count & 32 (5 safety, 14 core, 8 integration, 5 edge) \\
\bottomrule
\end{tabularx}
\end{center}

\begin{quotebox}
\textit{``The single greatest thing any individual can do for pollinators is to plant native wildflowers.''} --- Xerces Society

\vspace{0.5em}
Every flower is a node. Every garden is a hub. The network grows one planting at a time, requiring no central coordination, no permission, no infrastructure --- just soil, water, sun, and the decision to feed the bees. This is the architecture of resilience: distributed, autonomous, locally adapted, globally connected.
\end{quotebox}

\end{document}
